\documentclass[11pt]{article}

\usepackage[a4paper,margin=28mm]{geometry}
\usepackage{amsmath,amssymb,amsthm,mathtools}
\usepackage[hidelinks]{hyperref}
\usepackage{microtype}

\newtheorem{lemma}{Lemma}
\newtheorem{definition}{Definition}
\newtheorem{remark}{Remark}
\newcommand{\Z}{\mathbb Z}
\newcommand{\Rel}{\mathop{\rm Rel}}
\newcommand{\Sym}{\mathop{\rm Sym}}
\newcommand{\im}{\mathop{\rm im}}
\newcommand{\coker}{\mathop{\rm coker}}
\newcommand{\src}{\mathop{\rm src}}
\newcommand{\Obs}{\mathop{\rm Obs}}

\title{A Standalone Form of the Step--7 Obstruction:\\
  Terminal Lifting of a Signed Relation-Word Family}
\author{}
\date{}

\begin{document}
\maketitle

\section{Purpose}

This note removes almost all of the graded, metabelian, Smith, and
derived-functor notation from the missing Step--7 assertion.  It identifies a
single obstruction class.  Its vanishing is \emph{equivalent}, not merely
morally related, to the existence of the chain required by the completed
formal development.

The formulation also records an important negative fact.  The desired
primitive statement is not a general property of the ideal generated by the
relations in an enveloping algebra.  A proof must use the balance of the
complete signed occurrence family.  This is precisely why collecting one
relation-left word at a time leaves an unwanted component read.

\section{Terminal lifting data}

Let $F$ be a Lie ring whose underlying abelian group is finite free, and fix
an ordered integral PBW basis.  Write
\[
       P_j:U(F)\longrightarrow\Sym^jF
       \qquad (j=1,2)
\]
for the corresponding complete factor-$j$ projections, with
$\Sym^1F=F$.  Let $R\triangleleft F$ be a Lie ideal.

Let $\pi:F\twoheadrightarrow A$ be a surjective homomorphism to a
finite-free Lie ring and put $D=\pi(R)$.  Choose additive sections
\[
       \lambda:D\longrightarrow R,
       \qquad s:A\longrightarrow F
\]
of $R\twoheadrightarrow D$ and $\pi$, respectively.  Such sections exist in
the application because $A$ and $D$ are finite-free abelian groups.  Define
\[
 K=D\otimes A,
 \qquad B=\Sym^2A,
 \qquad
 \partial(d\otimes a)=d\odot a,
\]
and define the complete source-placement read by
\[
 \sigma:K\longrightarrow F,
 \qquad
 \sigma(d\otimes a)
   =P_1\bigl(\iota(\lambda d)\iota(sa)\bigr).
\]
Both maps are additive.

A finite signed full-relation word family is data
\[
       \mathcal W=\{(m_e,\rho_e,X_e)\}_{e\in E},
       \qquad
       m_e\in\Z,\quad \rho_e\in R,\quad X_e\in U(F),
\]
where each $\rho_e$ is retained as a whole element of $R$.  Set
\begin{align*}
 \Theta_{\mathcal W}&=\sum_{e\in E}m_e\,\iota(\rho_e)X_e,\\
 b_{\mathcal W}&=\Sym^2(\pi)P_2(\Theta_{\mathcal W})\in B,\\
 p_{\mathcal W}&=P_1(\Theta_{\mathcal W})\in F.
\end{align*}

\section{The smallest exact obstruction}

Consider the additive map
\[
 \Lambda:K\oplus R\longrightarrow B\oplus F,
 \qquad
 \Lambda(z,r)=\bigl(\partial z,\,\sigma z-r\bigr),
\]
where $R$ is included in $F$, and put
\[
       \mathcal O(F,R,\pi;\lambda,s)=\coker(\Lambda).
\]

\begin{definition}[terminal PBW obstruction]
The obstruction of $\mathcal W$ is
\[
       \Obs(\mathcal W)
       =\bigl[(b_{\mathcal W},p_{\mathcal W})\bigr]
       \in\mathcal O(F,R,\pi;\lambda,s).
\]
\end{definition}

\begin{lemma}[exact lifting criterion]\label{lem:criterion}
The following statements are equivalent.
\begin{enumerate}
\item $\Obs(\mathcal W)=0$.
\item There exist $z\in D\otimes A$ and $r\in R$ such that
\begin{equation}
       \partial z=b_{\mathcal W},
       \qquad
       \sigma z=p_{\mathcal W}+r.                 \label{eq:lift}
\end{equation}
\end{enumerate}
\end{lemma}

\begin{proof}
The obstruction is zero exactly when
$(b_{\mathcal W},p_{\mathcal W})$ lies in the image of $\Lambda$.  Writing
out the two coordinates of that assertion gives \eqref{eq:lift}.
\end{proof}

Thus the entire missing argument can be stated in one line.

\begin{lemma}[the missing exceptional vanishing]\label{lem:missing}
Let $\mathcal W_{\rm exc}$ be the signed family of full relation-left words
on the exceptional upper edge of the complete marked PBW/truncation trace.
Then
\[
                         \boxed{\Obs(\mathcal W_{\rm exc})=0.}
\]
\end{lemma}

This statement has no hidden choice of two unrelated chains.  The maps
$\partial$ and $\sigma$ use the same element $z$.  By
Lemma~\ref{lem:criterion}, a proof of Lemma~\ref{lem:missing} supplies exactly
the terminal chain and complete primitive equation needed by Step~7.

In the cyclic-top formulation one may replace $r\in R$ in
\eqref{eq:lift} by an error $e\in F$ whose evaluation belongs to the terminal
cyclic layer and has zero cyclic coordinate.  Injectivity of that coordinate
then implies that $e$ evaluates to zero, hence $e\in R$; the two formulations
are equivalent.

\section{Why ideal membership alone is insufficient}

Neither required PBW assertion is a formal consequence of
$\Theta_{\mathcal W}\in\iota(R)U(F)$.  The ordered PBW projection can split a
relation into basis coordinates before it records an interchange bracket.
Already with one spectator, the primitive assertion can fail; with two
spectators, the same phenomenon can obstruct the quadratic assertion.

For the primitive assertion, let $F$ be the free two-step nilpotent Lie ring
on ordered generators $a<b<c$, and let $R$ be the Lie ideal generated by
$a+c$.  With the PBW order beginning with $a<b<c$, take
\[
       \Theta=\iota(a+c)\iota(b).
\]
Since
\[
 \iota(c)\iota(b)=\iota(b)\iota(c)+\iota([c,b]),
\]
one has
\[
       P_1(\Theta)=[c,b].
\]
On the other hand, the degree-two part of $R$ is generated by
\[
       [c,a]
       \quad\text{and}\quad
       [a,b]+[c,b],
\]
so $[c,b]\notin R$.  Although
$P_2(\Theta)=(a+c)\odot b$ has a relation factor, the primitive equation
fails.

For the quadratic assertion, use the free two-step nilpotent Lie ring on
$a<b<c<d$, order the central commutator basis after the generators, keep the
same ideal $R=(a+c)$, and take
\[
       \Theta'=\iota(a+c)\iota(b)\iota(d).
\]
PBW collection gives
\[
       P_2(\Theta')=d\odot[c,b].
\]
This element is not in the image of
$R\otimes F\to\Sym^2F$.  Indeed, after tensoring with $\mathbb Q$ and passing
to $F/R$, its image is the nonzero product of the independent classes of
$d$ and $[c,b]$, whereas every symmetric product with an $R$-factor maps to
zero.  Thus even the existence of the desired factor-two chain is an
aggregate statement, not a wordwise ideal-membership lemma.

For the exceptional Step--7 aggregate, the missing partial brackets in this
example are replaced by the lower-context diagonal terms.  They disappear
only after the signed global trace is assembled and the horizontal and
vertical occurrences are paired.

\section{A detached certificate for proving the vanishing}

The following is a convenient finite certificate for
Lemma~\ref{lem:missing}.  It separates the genuinely difficult construction
from the final, elementary Stokes deduction.

Let $\overline F=F/R$, and let $C_1$ be the free abelian group on labelled
row occurrences and $C_2$ the free abelian group on labelled elementary
replacements.  Let
\[
       \delta:C_2\longrightarrow C_1
\]
send a replacement to its input occurrence minus its signed output
occurrences.  The literal enveloping-algebra identities define two additive
reads
\[
       \beta:C_1\longrightarrow B,
       \qquad
       \alpha:C_1\longrightarrow\overline F
\]
and give
\begin{equation}
                         \beta\delta=0,
       \qquad            \alpha\delta=0.             \label{eq:S0}
\end{equation}
The factor-two read is $\beta$; the map $\alpha$ is the complete factor-one
read followed by passage modulo the whole relation ideal.  Thus a full
relation is sent to zero only after all of its basis coordinates have been
grouped.

Call a finite ledger \emph{balanced for $\mathcal W$} if it contains an
element $S\in C_2$ and a decomposition
\begin{equation}
 \delta S=E-T+Q+
        \sum_{k}\bigl(V_k-H_k\bigr)                 \label{eq:S1}
\end{equation}
with the following properties:
\begin{center}
\begin{tabular}{c|c|c}
occurrence family & factor-two read & factor-one read in $\overline F$\\ \hline
$E$ (initial exceptional edge) & $b_{\mathcal W}$ & $[p_{\mathcal W}]$\\
$T$ (genuine terminal rows) & $\partial z$ & $[\sigma z]$\\
$Q$ (whole relation leaves) & $0$ & $0$\\
$H_k,V_k$ (paired diagonals) & $\beta(H_k)=\beta(V_k)$
  & $\alpha(H_k)=\alpha(V_k)$
\end{tabular}
\end{center}
for one and the same $z\in D\otimes A$.  Every equality in this table is an
equality of complete signed occurrence sums; no homogeneous component of a
relation is allowed in the $Q$ row.

\begin{lemma}[finite marked Stokes]\label{lem:stokes}
If $\mathcal W$ admits a balanced ledger, then
$\Obs(\mathcal W)=0$.
\end{lemma}

\begin{proof}
Apply $\beta$ to \eqref{eq:S1} and use \eqref{eq:S0}.  The diagonal terms
cancel and this gives $b_{\mathcal W}=\partial z$.  Apply $\alpha$ instead.
The diagonal terms again cancel, while the whole-relation leaves vanish, so
$[\sigma z]=[p_{\mathcal W}]$ in $F/R$.  Hence
$\sigma z=p_{\mathcal W}+q$ for some $q\in R$, and
Lemma~\ref{lem:criterion} applies.
\end{proof}

The nontrivial standalone task is therefore the following concrete existence
statement:
\[
 \boxed{\text{The exceptional signed relation-word family admits a balanced
 marked ledger.}}
\]
This is a finite combinatorial assertion.  Its construction must begin with
the complete occurrence trace, keep every relation whole, retain every
lower-context descendant, and form \eqref{eq:S1} only after internal
occurrences have cancelled.  In the application, the equality
$\alpha(H_k)=\alpha(V_k)$ is the detached form of
\[
       \Phi_k(f_k\chi_k)=T_k(\partial\chi_k),
\]
and the governing-coefficient calculation removes the remaining total
diagonal read.  The already proved nonexceptional and proper-subset families
may be deleted only after this global equality has been formed.

All objects in the certificate are finite free abelian groups, homomorphisms,
finite sums, and explicit equalities.  Hence both the obstruction formulation
and the balanced-ledger formulation can be stated directly in Lean.  Proving
the boxed existence statement supplies the missing chain without using the
circular ``known boundary chain minus subset tail'' construction.

\end{document}
